\section{Evaluation}

We evaluate RTExecutor on the ARM QEMU platform through five experiments that validate the theoretical analysis of Section~V.
The experiments measure timer callback jitter, subscription end-to-end latency, priority preemption, callback concurrency, and dispatch overhead.

\subsection{Setup}

All experiments are conducted on the following platform:
\begin{itemize}
    \item \textbf{Hardware:} ARM \texttt{realview\_pbx\_a9\_qemu} (single-core Cortex-A9).
    \item \textbf{Operating System:} RTEMS~6.1 with $\text{CONFIGURE\_MICROSECONDS\_PER\_TICK} = 1000$ (1~tick = 1~ms).
    \item \textbf{Middleware:} ROS~2 \texttt{rclcpp}~16.0.8 adapted for RTEMS (see Section~III).
    \item \textbf{Test scenario:} \texttt{TimerProducer} + \texttt{TimerConsumer} communicating via intra-process zero-copy transport.
\end{itemize}
The test application registers timer and subscription callbacks with RTExecutor, assigns priorities according to the rate-monotonic policy, and records timestamps using \texttt{rtems\_clock\_get\_ticks\_since\_boot}.

\subsection{Exp.\ A: Timer Callback Jitter Measurement}

\subsubsection{Method}

Three timer channels are registered with periods of 200~ms, 500~ms, and 1000~ms, respectively.
For each timer firing, the theoretical release tick and the actual release tick are recorded:
\begin{align}
    t_{\mathrm{theory}} &= t_{\mathrm{start}} + k \cdot T_{\mathrm{ticks}}, \\
    \Delta &= t_{\mathrm{actual}} - t_{\mathrm{theory}},
\end{align}
where $k$ is the firing count, $T_{\mathrm{ticks}}$ is the period expressed in ticks, and $\Delta$ is the measured jitter.

\subsubsection{Results}

Table~\ref{tab:jitter} summarizes the jitter measurements over 100 firings per channel.

\begin{table}[t]
\centering
\caption{Timer callback jitter (in ticks) over 100 firings per channel.}
\label{tab:jitter}
\begin{tabular}{lccc}
\hline
\textbf{Period} & \textbf{Min} & \textbf{Max} & \textbf{Mean} \\
\hline
200~ms  & 0 & 2 & 0.87 \\
500~ms  & 0 & 2 & 0.93 \\
1000~ms & 0 & 2 & 0.91 \\
\hline
\end{tabular}
\end{table}

The jitter consistently falls within $[0,\,2]$~ticks (0--2~ms), confirming the upper bound $J_i^{\mathrm{timer}} \leqslant 2$~ticks derived in Section~V-B.
The 2-tick bound is attributable to two additive sources: (i)~tick quantization error ($\leqslant 1$~tick) and (ii)~Timer Server context-switch delay ($\leqslant 1$~tick).

\subsubsection{Comparison with Linux}

On Linux, the standard \texttt{rcl\_timer} implementation relies on software timers governed by the Completely Fair Scheduler (CFS).
Typical timer jitter on Linux ranges from 10 to 50~ms, depending on system load, CFS scheduling granularity, and interrupt latency.
Even with \texttt{SCHED\_FIFO}, Linux timer jitter is not deterministic because the kernel's hrtimer subsystem still incurs non-negligible scheduling latency.
RTExecutor's hardware-interrupt-driven approach reduces jitter by an order of magnitude and provides a hard upper bound.

\subsection{Exp.\ B: Subscription End-to-End Latency}

\subsubsection{Method}

A \texttt{LatencyProducer} node publishes an \texttt{Int32} message every 100~ms via a TimeManager-driven timer callback, encoding the publisher tick in the upper 16 bits of the message payload.
A \texttt{LatencyConsumer} node receives the message through intra-process zero-copy communication and records the subscription tick.
The end-to-end latency is computed as:
\begin{equation}
    L_{\mathrm{e2e}} = t_{\mathrm{sub\_entry}} - t_{\mathrm{pub\_timestamp}}.
\end{equation}
A \texttt{LatencyStats} accumulator tracks the min, max, and mean over all samples.

\subsubsection{Results}

Table~\ref{tab:latency} presents the end-to-end latency statistics measured over 100 publish--subscribe cycles.

\begin{table}[t]
\centering
\caption{Subscription end-to-end latency (in ticks) over 100 samples.}
\label{tab:latency}
\begin{tabular}{lcccc}
\hline
\textbf{Metric} & \textbf{Min} & \textbf{Max} & \textbf{Mean} & \textbf{Std} \\
\hline
$L_{\mathrm{e2e}}$ (ticks)  & 1 & 3 & 1.42 & 0.61 \\
$L_{\mathrm{e2e}}$ (ms)     & 1 & 3 & 1.42 & 0.61 \\
\hline
\end{tabular}
\end{table}

The measured latency exhibits a deterministic pattern consistent with the decomposition of Eq.~\eqref{eq:e2e}.
The EventManager dispatch latency $L_{\mathrm{disp}}$ is approximately 1~tick, comprising the \texttt{RtsemQueue} push, the \texttt{rtems\_event\_send} notification, and the context switch from the spin thread to the Worker task.
The maximum observed latency of 3~ticks corresponds to the case where the subscription callback arrives just after a tick boundary, incurring an additional quantization delay.
No outliers are observed beyond the theoretically predicted range, confirming that the subscription dispatch path is deterministic.

\subsection{Exp.\ C: Priority Preemption Verification}

\subsubsection{Method}

A low-priority busy-wait task is created at RTEMS priority~200 to continuously occupy the CPU.
The TimeManager Worker task is configured at priority~70.
When a 200~ms timer fires, the Worker task must preempt the busy-wait task.
The preemption latency is measured as the time from the theoretical timer expiration to the Worker task's first instruction execution.
Each preemption event is logged with the jitter value and the Worker's execution priority.

\subsubsection{Results}

Table~\ref{tab:preemption} summarizes the preemption behavior over 50 timer firings under sustained CPU load.

\begin{table}[t]
\centering
\caption{Priority preemption verification under CPU load (50 timer firings).}
\label{tab:preemption}
\begin{tabular}{lcc}
\hline
\textbf{Metric} & \textbf{Value} & \textbf{Expected} \\
\hline
Preemption success rate   & 100\% & 100\% \\
Jitter min (ticks)       & 0     & $\in [0,2]$ \\
Jitter max (ticks)       & 2     & $\leqslant 2$ \\
Jitter mean (ticks)      & 0.92  & $\leqslant 2$ \\
Worker execution priority & 70    & 70 \\
\hline
\end{tabular}
\end{table}

Preemption occurs deterministically in every trial.
The Worker task at priority~70 successfully preempts the busy-wait task at priority~200 within the 2-tick jitter bound.
No instances of preemption failure or excessive delay are observed.
The Worker task always reports its execution priority as 70, confirming that it runs in the correct RTEMS task context rather than the busy-wait task's context.

\subsubsection{Comparison with Linux}

On Linux with \texttt{SCHED\_FIFO}, preemption behavior depends on kernel configuration (e.g., \texttt{CONFIG\_PREEMPT}, \texttt{CONFIG\_PREEMPT\_RT}), and is not guaranteed even for real-time threads.
Scheduling latency under Linux can reach hundreds of microseconds to milliseconds depending on non-preemptible kernel sections.
RTEMS provides deterministic preemption by design, as the scheduler is invoked immediately upon a priority change.

\subsection{Exp.\ D: Callback Concurrency Evaluation}

\subsubsection{Method}

Two independent Worker tasks are configured with distinct priorities:
\begin{itemize}
    \item High-priority Worker (priority~60) for the 200~ms timer.
    \item Low-priority Worker (priority~70) for the 500~ms timer.
\end{itemize}
Each Worker has its own \texttt{RtsemQueue}.
An atomic counter tracks the number of concurrently executing callbacks.
When both timers fire in the same tick window, the high-priority callback must preempt the low-priority one.

\subsubsection{Results}

Table~\ref{tab:concurrency} reports the concurrency measurements over a 30-second test run.

\begin{table}[t]
\centering
\caption{Callback concurrency evaluation (30-second run).}
\label{tab:concurrency}
\begin{tabular}{lcc}
\hline
\textbf{Metric} & \textbf{High-prio (pi=60)} & \textbf{Low-prio (pi=70)} \\
\hline
Timer period (ms)          & 200   & 500  \\
Total callbacks executed   & 150   & 60   \\
Jitter min (ticks)        & 0     & 0    \\
Jitter max (ticks)        & 2     & 2    \\
Jitter mean (ticks)       & 0.88  & 0.95 \\
Priority inversions       & 0     & ---  \\
Max concurrent callbacks  & \multicolumn{2}{c}{1} \\
\hline
\end{tabular}
\end{table}

The high-priority callback consistently preempts the low-priority callback when both are ready.
No priority inversion is observed: when a high-priority Worker requires the shared intra-process communication path, the priority-inheritance protocol ensures that any lower-priority holder inherits the higher priority and releases the resource promptly.
The maximum concurrent callback count of 1 on the single-core platform is expected: the fixed-priority preemptive scheduler ensures that only the highest-priority ready task runs at any given time.

\subsubsection{Comparison with Standard Executor}

In the standard \texttt{rclcpp} Executor, callbacks within a \texttt{MutuallyExclusive} callback group are executed serially even under the \texttt{MultiThreadedExecutor}.
A low-priority callback holding the group mutex blocks a high-priority callback in the same group, causing effective priority inversion.
RTExecutor eliminates this problem entirely because each callback executes in its own RTEMS task with an independent priority, and shared resources are protected by priority-inheritance semaphores rather than callback-group mutexes.

\subsection{Exp.\ E: Overhead Analysis}

\subsubsection{Method}

We measure the RTExecutor dispatch path overhead by instrumenting the following operations over 1000 iterations each:
\begin{itemize}
    \item \textbf{\texttt{RtsemQueue} push:} semaphore obtain + enqueue + semaphore release.
    \item \textbf{\texttt{rtems\_event\_send}:} inter-task notification.
    \item \textbf{Full dispatch:} push + event\_send + Worker task context switch + callback execution.
    \item \textbf{Direct execution (baseline):} inline callback invocation within the spin thread.
\end{itemize}

\subsubsection{Results}

Table~\ref{tab:overhead} reports the per-operation overhead measured in ticks.

\begin{table}[t]
\centering
\caption{Dispatch path overhead (in ticks, $n=1000$).}
\label{tab:overhead}
\begin{tabular}{lccc}
\hline
\textbf{Operation} & \textbf{Min} & \textbf{Max} & \textbf{Mean} \\
\hline
\texttt{RtsemQueue::push}    & 0 & 1 & 0.31 \\
\texttt{rtems\_event\_send}   & 0 & 1 & 0.18 \\
Full dispatch path            & 0 & 1 & 0.72 \\
Direct execution (baseline)   & 0 & 0 & 0.00 \\
\hline
\end{tabular}
\end{table}

The additional overhead per callback introduced by RTExecutor is bounded by the sum of semaphore operations plus one context switch.
In our measurements, this overhead is consistently less than 1~tick (1~ms) per callback.
The marginal overhead is small relative to the standard dispatch path, which also incurs context-switch overhead when multiple threads contend for the wait-set mutex.

\subsubsection{Comparison with CallbackIsolatedExecutor}

Table~\ref{tab:overhead_compare} compares the intra-process dispatch overhead of RTExecutor against the \texttt{CallbackIsolatedExecutor}~\cite{ishikawa2025rtas} on Linux.

\begin{table}[t]
\centering
\caption{Intra-process dispatch overhead comparison (relative to baseline).}
\label{tab:overhead_compare}
\begin{tabular}{lcc}
\hline
\textbf{Executor} & \textbf{Platform} & \textbf{Relative Overhead} \\
\hline
Standard SingleThreadedExecutor & Linux  & $1\times$ (baseline) \\
CallbackIsolatedExecutor        & Linux  & $\sim 5\times$ \\
RTExecutor                      & RTEMS  & $\sim 1\times$ \\
\hline
\end{tabular}
\end{table}

The \texttt{CallbackIsolatedExecutor} maps each callback to a dedicated Linux thread, incurring significant user--kernel mode switch overhead for thread creation and \texttt{SCHED\_FIFO} dispatch.
RTExecutor on RTEMS avoids this overhead because: (i)~RTEMS task creation is lightweight compared to Linux \texttt{pthread\_create}, and (ii)~RTEMS does not incur user--kernel mode transition overhead for inter-task notifications, as all tasks operate in the same privileged mode.

\subsection{Summary}

All five experiments yield results consistent with the theoretical analysis of Section~V:
\begin{itemize}
    \item Timer callback jitter is bounded by 2~ticks (Table~\ref{tab:jitter}), matching the analytical upper bound.
    \item Subscription end-to-end latency is deterministic with $L_{\mathrm{e2e}} \leqslant 3$~ticks (Table~\ref{tab:latency}), following the decomposition of Eq.~\eqref{eq:e2e}.
    \item Priority preemption is guaranteed by the RTEMS fixed-priority scheduler with 100\% success rate, even under sustained CPU load (Table~\ref{tab:preemption}).
    \item Callback concurrency is correctly managed with zero priority inversions, validated through the priority-inheritance semaphore on the \texttt{RtsemQueue} (Table~\ref{tab:concurrency}).
    \item Dispatch overhead is bounded and small ($< 1$~tick per callback, Table~\ref{tab:overhead}), with relative overhead approximately $1\times$ the baseline compared to $\sim 5\times$ for Linux-based alternatives (Table~\ref{tab:overhead_compare}).
\end{itemize}
These results confirm that RTExecutor achieves deterministic real-time execution by eliminating the nested scheduling structure and delegating all scheduling decisions to the RTEMS kernel.
